\section{An exact reciprocal-growth criterion for the unrestricted one-sided operator}
\label{sec:reciprocal-growth-criterion}

This section concerns the actual one-sided Hilbert space and
compact-frequency core of MA-RS-001--007.  Its universal quantifier
is over smooth physical cutoffs near one fixed cutoff.  It does
not decide a prescribed individual cutoff, and it does not
assert that the arithmetic reciprocal has nonpolynomial growth.
The completed-ideal return of MA-MO-028 has a different domain
and remains unconditional.

Fix a finite set \(S\) of rational primes and
\(\sigma\in\mathscr G\) as in (MA-RS.2).  Retain
\[
\begin{gathered}
 a=\sigma-\tfrac12,\qquad
 \mathcal H_a=L^2(\mathbb R,e^{2ax}\,dx),\qquad
 V_a=\mathcal F_+M_{e^{ax}},\\
 \mathcal F_+f(t)=(2\pi)^{-1/2}\int e^{itx}f(x)\,dx,\qquad
 \mathscr C=C_c^\infty(\mathbb R),\qquad
 \mathscr D_a=V_a^{-1}\mathscr C,\\
 A(t)=\zeta_S(\sigma+it),\qquad j(t)=A(t)^{-1}.
\end{gathered}
\tag{MA-MO.1000}
\]
All derivatives of \(A\) have polynomial growth.  The
already proved generic-line estimates give
\[
 |j^{(k)}(t)|\le
       \exp(C_k\log^2(2+|t|)),\qquad
 \left|\frac{j'(t)}{j(t)}\right|
                       \le C(2+|t|)^D .
\tag{MA-MO.1001}
\]
The second assertion is (MA-RS.64), with the derivative
converted to the real frequency variable.  In particular
\(A\) and \(j\) are smooth and nonzero on the real line.
The maximal domain of \(M_j\) is
\(\{v\in L^2:jv\in L^2\}\).

Fix a nonzero, possibly complex, test
\(g\in C_c^\infty(\mathbb R)\), and write
\[
 G(z)=\int e^{zx}g(x)\,dx,\qquad m(t)=G(a+it).
\tag{MA-MO.1002}
\]
There is no restriction on \(g(0)\); in particular the
theorem applies to every nonzero test with \(g(0)=0\).
Choose a real smooth step \(\phi_0\), equal to one
far left and zero far right, such that
\(0\le\phi_0\le1\) and \(\phi_0'\) has compact support
and finite Gevrey order.  Choose a closed bounded
interval \(I\) with nonempty interior such that
\[
 0<\delta\le\phi_0(x)\le1-\delta\quad(x\in I).
\]
Let \(\mathscr D_I^\mathbb R\) be the real Frechet
space of real smooth functions supported in \(I\), with
seminorms \(\max_{0\le k\le N}\|\psi^{(k)}\|_\infty\).
Fix \(0<\epsilon<\delta\) and put
\[
 \mathscr U=\{\psi\in\mathscr D_I^\mathbb R:
                         \|\psi\|_\infty<\epsilon\},\qquad
 \phi_\psi=\phi_0+\psi,\qquad
 T_\phi=\mathcal F_+M_\phi\mathcal F_+^{-1}.
\tag{MA-MO.1003}
\]
Every \(\phi_\psi\) is a real smooth step with
\(0\le\phi_\psi\le1\) and compactly supported derivative.
Monotonicity and reflection symmetry are not imposed.

Since \(g\ne0\), the entire transform \(G\) is not
identically zero.  Choose a bounded open frequency
interval \(\Omega\) on which \(mA\ne0\), and fix once
and for all
\[
 0\ne q\in C_c^\infty(\Omega),\qquad
 w=\frac{q}{\overline{mA}}\in C_c^\infty(\Omega),\qquad
 f=\mathcal F_+^{-1}q .
\tag{MA-MO.1004}
\]
The vector \(q\) and the output-testing vector \(w\)
do not vary with the cutoff perturbation.  The function
\(f\) is entire and is not identically zero.

Throughout this section a step \(\phi\) has the normalized
endpoints: it is identically one far left and identically
zero far right.  For every such smooth bounded step with
\(\phi'\in C_c^\infty\), define on the same fixed core
\(\mathscr C\) the actual operator
\[
 K^\circ_{g,\phi}v
      =M_mT_\phi v-M_{mA}T_\phi M_jv,\qquad v\in\mathscr C .
\tag{MA-MO.1005}
\]
Both terms in this formula are defined and are Schwartz
functions: \(jv\) is compact smooth, \(T_\phi\) maps
Schwartz functions to Schwartz functions, and \(A\)
has polynomially bounded derivatives.  Thus (MA-MO.1005)
does not use a hypothetical bounded inverse or an assigned
trace.  On any vector where the sequential commutator
is defined, it agrees with
\(V_aC_gZ_{a,S}[J_{a,S},M_\phi]V_a^{-1}\).
It need not imply that this latter sequential expression
is defined on every vector of \(\mathscr C\).

\paragraph{MA-MO-031 (polynomial growth, universal cutoff scope, and actual failure witnesses).}
In the fixed setting above, the following conditions
are equivalent:
\[
\begin{array}{ll}
\mathrm{(P)}&
 |j(t)|\le C(2+|t|)^N
          \text{ for some }C,N\text{ and every }t;\\[2pt]
\mathrm{(D)}&
 T_{\phi_\psi}q\in\mathcal D(M_j)
                           \text{ for every }\psi\in\mathscr U;\\[2pt]
\mathrm{(B)}&
 K^\circ_{g,\phi_\psi}\text{ has a bounded extension to }L^2
                           \text{ for every }\psi\in\mathscr U;\\[2pt]
\mathrm{(T)}&
 K^\circ_{g,\phi_\psi}\text{ has a trace-class extension to }L^2
                           \text{ for every }\psi\in\mathscr U .
\end{array}
\tag{MA-MO.1006}
\]
In fact (P) implies a stronger positive conclusion:
for every compact smooth test \(h\) and every smooth
bounded step \(\phi\) with compactly supported derivative,
the original sequential commutator is defined on
\(\mathscr D_a\), has the Schwartz kernel (MA-RS.16),
and extends to a trace-class operator.  The actual
finite-frequency operators (MA-RS.54) converge to it
in trace norm.  Its trace and residue correction are
exactly (MA-RS.56).

If (P) fails, the following set is a dense
\(G_\delta\) subset of \(\mathscr D_I^\mathbb R\):
\[
 \{\psi:
          jT_{\phi_0+\psi}q\notin L^2\}.
\tag{MA-MO.1007}
\]
Consequently every neighborhood of zero contains a
real perturbation \(\psi\in\mathscr U\) for which
the original sequential commutator fails its inverse
domain already on the fixed core vector \(q\).
For the same cutoff the always-defined operator
(MA-MO.1005) has no bounded extension.  There are
explicit normalized vectors \(v_R\in\mathscr C\)
with
\[
       \|K^\circ_{g,\phi_\psi}v_R\|_2\longrightarrow\infty.
\tag{MA-MO.1008}
\]
For the actual operators \(K_N(g)\) of (MA-RS.54)
with this fixed cutoff,
\[
 \sup_N\|K_N(g)\|_{\mathcal H_a\to\mathcal H_a}=\infty .
\tag{MA-MO.1009}
\]
Their traces nevertheless converge to the value
in (MA-RS.56).  In particular there is no bounded
strong operator limit and no trace-norm limit in
this conditional failure case.

\begin{proof}
\emph{1. Polynomial reciprocal growth gives all smooth
domains and the actual kernel.}
Assume (P).  Repeatedly differentiate \(Aj=1\).
Every derivative \(j^{(k)}\) is a finite sum of products
of powers of \(j\) and derivatives of \(A\).
It therefore has polynomial growth, with an exponent
allowed to depend on \(k\).  Thus both \(A\) and \(j\)
are continuous multipliers of the Schwartz space.
For a core vector \(v\), the functions
\(T_\phi v\) and \(T_\phi jv\) are Schwartz.
It follows that \(M_jT_\phi v\) is in the maximal
domain of \(M_A\), and that \(T_\phi jv\) is also in
that domain.  Moreover
\(M_AM_jT_\phi v=T_\phi v\).
This checks the two actual sequential compositions
and gives (MA-MO.1005) as their difference.

For \(h\in C_c^\infty\), put
\(m_h(t)=\mathcal Mh(a+it)\), and put
\(p(\xi)=\int e^{i\xi x}\phi'(x)\,dx\).
Both are Schwartz functions.  The exact distribution
calculation (MA-RS.59)--(MA-RS.60) gives the kernel
\[
 k_{h,\phi}(t,u)
  =\frac{i}{2\pi}m_h(t)A(t)
       \left(\int_0^1j'(u+\theta(t-u))\,d\theta\right)
                                      p(t-u).
\tag{MA-MO.1010}
\]
Every derivative of the parenthesized function has
polynomial growth in \((t,u)\).  The variables \(t\)
and \(t-u\) form invertible linear coordinates on
\(\mathbb R^2\).  Schwartz decay of \(m_h(t)\) and
\(p(t-u)\), with all their derivatives, consequently
proves \(k_{h,\phi}\in\mathcal S(\mathbb R^2)\).

Expansion in the two-variable Hermite basis has
coefficients decaying faster than every power of
both indices.  The resulting absolutely summable
rank-one series proves trace class, and its trace
is the integral of the diagonal.  Equivalently this
is the Schwartz-kernel argument already used in
MA-RS-002/004.  This implication is continuous
from Schwartz kernels to trace class, since finitely
many oscillator seminorms bound the sum of the
absolute coefficients.

Let \(\chi_N(t)=\chi(t/N)\) be exactly the cutoffs
in MA-RS-007.  Formula (MA-RS.54) agrees with
the kernel \(\chi_N(t)k_{h,\phi}(t,u)\chi_N(u)\).
Since \(\chi=1\) near zero, this converges to
\(k_{h,\phi}\) in every Schwartz seminorm.
The preceding nuclear bound proves trace-norm
convergence.  The diagonal, its sign and its
residue evaluation are therefore exactly
(MA-RS.56), now as a limit of operators as well
as of traces.  This proves (P) implies (D) and (T);
(T) implies (B) directly.

\emph{2. A fixed output functional forces the fixed
inverse-domain condition.}
The step is real, so \(T_\phi\) is self-adjoint.
Use the Hilbert inner product linear in its first
variable.  For \(v\in\mathscr C\), the actual two
compositions give
\[
 \begin{split}
 \langle K^\circ_{g,\phi}v,w\rangle
 &=\left\langle v,\,
       T_\phi(\overline m\,w)
            -\overline j\,T_\phi q\right\rangle,\\
 H_\phi&=T_\phi(\overline m\,w)
            -\overline j\,T_\phi q .
 \end{split}
\tag{MA-MO.1011}
\]
The pairing on the right is well defined as an
ordinary integral over the compact support of \(v\).
Indeed \(H_\phi\) is a smooth function, though it
need not be in \(L^2\).  The identity uses precisely
\(\overline{mA}\,w=q\), with no restriction on the
phase of the test \(g\).

If \(K^\circ_{g,\phi}\) has a bounded extension,
the left side is a bounded functional of \(v\)
in \(L^2\).  The Riesz representation and equality
on all compact smooth vectors imply that the
smooth function \(H_\phi\) is in \(L^2\).
Its first term is already in \(L^2\), since
\(T_\phi\) is bounded.  Hence
\(\overline j\,T_\phi q\in L^2\), or equivalently
\(jT_\phi q\in L^2\).  This proves (B) implies (D).

\emph{3. Universal cutoff domains force a polynomial
local reciprocal bound.}
The fixed Gevrey step satisfies
\[
                 jT_{\phi_0}q\in L^2 .
\tag{MA-MO.1012}
\]
For completeness, outside a fixed interval containing
the support of \(q\), the transform \(T_{\phi_0}q\)
is the ordinary integral (MA-RS.20).  The Fourier
transform of \(\phi_0'\) has
\(\exp(-c|t|^\theta)\) decay for some \(\theta>0\).
This beats the \(\exp(C\log^2(2+|t|))\) growth
of \(j\), and proves (MA-MO.1012).

Assume (D).  Subtract (MA-MO.1012) and scale each
real perturbation by a sufficiently small nonzero
constant to put it in \(\mathscr U\).  One obtains
\[
 \mathcal L\psi
       :=j\,\mathcal F_+(\psi f)\in L^2
                  \quad(\psi\in\mathscr D_I^\mathbb R).
\tag{MA-MO.1013}
\]
Taking real and imaginary parts extends this
assertion to the complex Frechet space \(\mathscr D_I\).
The linear map \(\mathcal L:\mathscr D_I\to L^2\)
is continuous.  Here is an explicit Baire argument.
For each positive integer \(n\), the set
\(\{\psi:\|\mathcal L\psi\|_2\le n\}\) is closed:
convergence in \(\mathscr D_I\) gives pointwise
convergence of the Fourier transforms, and Fatou's
lemma gives the norm inequality.  These closed
sets cover \(\mathscr D_I\).  One has interior,
so taking differences gives a neighborhood of
zero on which the norm is bounded.  Homogeneity
therefore yields finite constants \(C,N\) with
\[
 \|j\mathcal F_+(\psi f)\|_2
       \le C\max_{0\le k\le N}\|\psi^{(k)}\|_\infty
                         \quad(\psi\in\mathscr D_I).
\tag{MA-MO.1014}
\]

The nonzero entire function \(f\) has only isolated
real zeros.  Choose a closed subinterval \(J\)
in the interior of \(I\) on which \(f\ne0\).
Take a nonnegative nonzero \(b\in C_c^\infty(J)\).
Its normalized Fourier transform is nonzero at
zero; for some \(\delta_1,c_1>0\),
\[
 |\mathcal F_+b(t)|\ge c_1\quad(|t|\le\delta_1).
\]
For each real \(T\), the compact smooth function
\[
          \psi_T(x)=\frac{b(x)}{f(x)}e^{-iTx}
\tag{MA-MO.1015}
\]
belongs to \(\mathscr D_I\), and its derivative
seminorm through order \(N\) is
\(O((2+|T|)^N)\).  Substitution in
(MA-MO.1014) gives
\[
 \int_{T-\delta_1}^{T+\delta_1}|j(t)|^2\,dt
                        \le C'(2+|T|)^{2N}.
\tag{MA-MO.1016}
\]
The use of complex \(\psi_T\) here is legitimate
because the real domain assertion was first extended
complex linearly; the failure cutoff constructed
below remains real.

\emph{4. Local \(L^2\) control implies pointwise
polynomial growth on the generic line.}
By (MA-MO.1001),
\[
 \left|\frac{d}{dt}\log|j(t)|\right|
                        \le C(2+|t|)^D.
\]
Choose
\[
 r_T=\min\{\delta_1/2,\ 1/2,\
                       [4C(3+|T|)^D]^{-1}\}.
\tag{MA-MO.1017}
\]
Enlarging \(C\) if necessary, integration on
\([T-r_T,T+r_T]\) proves
\(|j(t)|\ge e^{-1/2}|j(T)|\) there.
Combining with (MA-MO.1016) yields
\[
 2e^{-1}r_T|j(T)|^2
       \le C'(2+|T|)^{2N}.
\tag{MA-MO.1018}
\]
Since \(r_T\ge c(2+|T|)^{-D}\), this is a
polynomial pointwise bound for \(j(T)\).
Thus (D) implies (P), proving all equivalences.
The polynomial logarithmic-derivative estimate is
used exactly here; a smooth zero-free line without
that estimate is not substituted for the generic
hypothesis.

\emph{5. Real smooth failure cutoffs are residual.}
Suppose (P) fails.  On the real space
\(\mathscr D_I^\mathbb R\), define
\[
 E_n=\{\psi:
       \|j\mathcal F_+(\psi f)\|_2\le n\},
\tag{MA-MO.1019}
\]
where a nonintegrable square norm is assigned
\(+\infty\).  Each \(E_n\) is closed by the
same pointwise convergence and Fatou argument.
If any \(E_n\) had interior, its difference set
would imply that every sufficiently small real
perturbation has finite square norm.  Scaling
would then give that property for every real
perturbation, and complex linearity and
steps3--4 would prove (P), a contradiction.
Thus each \(E_n\) has empty interior.
Their union is exactly the inverse-domain set
for the perturbation part.  Its complement is
a dense \(G_\delta\), by Baire's theorem.
Adding the already integrable base term
(MA-MO.1012) leaves membership unchanged.
This proves (MA-MO.1007).

Choose a bad \(\psi\) in \(\mathscr U\).
The real cutoff retains all plateau and range
conditions in (MA-MO.1003), but \(M_jT_\phi q\)
does not exist in \(L^2\).  Thus the original
sequential expression is not defined on the
entire fixed core.  At the same time
the actual function \(H_\phi\) in
(MA-MO.1011) is not in \(L^2\).
Choose real \(\eta\in C_c^\infty\) with
\(0\le\eta\le1\), equal to one near zero,
and put \(\eta_R(t)=\eta(t/R)\).  For all
sufficiently large \(R\), define
\[
 v_R=\frac{\eta_R H_\phi}{\|\eta_R H_\phi\|_2}
                                      \in\mathscr C.
\tag{MA-MO.1020}
\]
The denominator tends to infinity: for each
fixed compact interval the cutoff is eventually
one, and \(H_\phi\notin L^2\).
By (MA-MO.1011) and \(0\le\eta_R\le1\),
\[
 \begin{split}
 |\langle K^\circ_{g,\phi}v_R,w\rangle|
 &=\frac{\int\eta_R|H_\phi|^2}
              {\|\eta_R H_\phi\|_2}\\
 &\ge\|\eta_R H_\phi\|_2\longrightarrow\infty .
 \end{split}
\tag{MA-MO.1021}
\]
This proves the actual normalized-core divergence,
including for arbitrary complex nonzero \(g\).
The physical vectors \(V_a^{-1}v_R\) have norm
one in \(\mathcal H_a\).  No assertion that
they belong to the two-sided completed ideal
is made.

\emph{6. The finite operators cannot converge,
although their traces do.}
Every finite operator (MA-RS.54) exists and is
trace class for this same smooth cutoff.  For
a fixed \(v\in\mathscr C\), \(\chi_Nv=v\)
for all sufficiently large \(N\); the same
is true of \(jv\).  Consequently the exact
bounded-composition formula gives
\[
 V_aK_N(g)V_a^{-1}v
        =\chi_NK^\circ_{g,\phi}v
                 \longrightarrow K^\circ_{g,\phi}v
                                      \quad\hbox{in }L^2 .
\tag{MA-MO.1022}
\]
If the operator norms of \(K_N(g)\) were
uniformly bounded, this identity on the
dense fixed core would give a bounded
extension of \(K^\circ_{g,\phi}\), contrary
to (MA-MO.1021).  Hence (MA-MO.1009) holds.
The uniform boundedness principle also rules
out a bounded strong operator limit.
Norm and trace-norm convergence are excluded
in particular.  The trace calculation of
MA-RS-007 uses only the polynomial logarithmic
derivative, not a polynomial reciprocal.
It therefore still gives the convergent
traces and their exact residue correction.
\end{proof}

\paragraph{A concrete sufficient zero configuration for the conditional failure branch.}
The analytic hypothesis that (P) fails can be implied
by the following explicit, but not asserted, arithmetic
configuration.  Fix \(R=1/4\).  Suppose there are
actual nontrivial zeros
\(\rho_n=\beta_n+i\gamma_n\), of multiplicity \(m_n\),
with \(|\gamma_n|\to\infty\), such that
\[
 \begin{gathered}
 0<d_n=|\sigma-\beta_n|<R/2,\\
 \frac{m_n\log(R/d_n)}
                   {\log(2+|\gamma_n|)}\longrightarrow\infty .
 \end{gathered}
\tag{MA-MO.1023}
\]
Then (P) fails and the real-cutoff/core/operator
conclusions above apply.

Indeed the fixed-strip upper bound gives
\(|\zeta(s)|\le C(2+|\gamma_n|)^C\) on the
circle \(|s-\rho_n|=R\), for sufficiently large
\(n\).  The disk contains no pole of zeta.
The function \(\zeta(s)/(s-\rho_n)^{m_n}\)
is holomorphic throughout it, regardless of
other zeros in the disk.  The maximum principle
therefore gives
\[
 |\zeta(\sigma+i\gamma_n)|
        \le C(2+|\gamma_n|)^C(d_n/R)^{m_n}.
\tag{MA-MO.1024}
\]
The finitely many Euler factors \(E_S\) are
uniformly bounded and bounded away from zero
on the fixed line \(\sigma>0\).  Thus
\[
 |j(\gamma_n)|
       \ge c(2+|\gamma_n|)^{-C}(R/d_n)^{m_n},
\tag{MA-MO.1025}
\]
which exceeds every fixed polynomial along
this sequence under (MA-MO.1023).
This is a conditional implication, not evidence
that such zeros or multiplicities occur.
The generic-distance bound (MA-RS.2) alone
does not establish (MA-MO.1023), and no result
in this section asserts nonpolynomial growth
of an actual arithmetic reciprocal.

\paragraph{Source, completion and remaining scope.}
The source for the actual closed multipliers,
sequential domains and Fourier sign is
MA-RS.8--9 and MA-RS.15--21; the finite
operators and polynomial logarithmic-derivative
input are MA-RS.52--64.  Their finite-disk
and selected-height input is the supplied
MT-M18 proof, (M18.7)--(M18.9), already
audited in those statements.  The positive
Schwartz-kernel argument also recovers
MA-RS-004 when its fixed zero gap supplies
(P).  No new universal zero estimate is
imported here.

In the rational setting of this section, the
user-stipulated input \(A\) supplies the fixed
gap \(\sigma-1/2>0\), and MA-RS-004 therefore
proves (P) and all the positive operator
conclusions under \(A\).  The unresolved
arithmetic issue below is an independent
proof without that zero-location input.
Using this \(A\)-supported implication to
justify an independent narrow trace argument
would retain the same arithmetic premise.

The theorem sharpens the unrestricted operator
gate in \texttt{MA-Q1-RF002}/\texttt{MA-Q1-RS003}
to a precise analytic condition for the stated
universal cutoff class.  It does not independently
decide that arithmetic condition without
zero-location input, an individual
predetermined cutoff, or singular boundary
lines.  Its normalized countervectors are
vectors of the one-sided Hilbert core; they
are not asserted to lie in the actual
two-sided completed theta ideal.  In
MA-MO-028 that ideal cancels the full divisor
before contour deformation, so its
unconditional return is compatible with
every assertion here.  The polynomial
partial-return costs of MA-PR and the
full-source divisor premises of
MA-MO-027/030 remain different questions.

In the inherited Grothendieck trace framework,
the relevant additional structure is the
actual continuous nuclear map and its
domain, not a scalar assignment of its
trace.  The arithmetic Euler determinants
and finite-length divisor factors are
unchanged.  The trace/determinant programme
recorded in cote~1 BR-TR-01/02, scans~079--080
and~111--112, supplies no implication from
convergent scalar traces to bounded
operator convergence; (MA-MO.1009) and
(MA-RS.56) explicitly separate those
properties in the conditional failure
case.  These archive locators are inherited,
not claimed freshly read for this theorem.
